% !TEX root = final_report.tex
\documentclass[final_report.tex]{subfiles}

% Appendices are not counted towards the 60-page content limit, and are not
% guaranteed to be read. Nothing load-bearing belongs here --- and no source
% code listings, which live in the repository.

\begin{document}
\apxsection{Implementation reference}\label{apx:implementation-tables}
\Cref{sec:implementation} describes the test suite and the two accelerators the code was run
on; the two tables below are the reference detail behind those descriptions.

\begin{table}[htbp]
  \centering
  \caption[What the test suite guarantees]{What the test suite guarantees. Each row pairs a claim
  made in \cref{ch:ropart} with the test that holds it; these ten are the
  chapter-relevant subset of the 37-test suite, the remainder covering shapes,
  registries and the evaluation utilities. Note that they are property and oracle checks, rather than
  unit tests. The suite is run locally under MPS, it has not been run under CUDA.}
  \label{tab:tests}
  \footnotesize
  \begin{tabular}{p{0.42\textwidth} l}
    \toprule
    Claim & Test \\
    \midrule
    Targets equal an independent oracle (\cref{eq:target-ropart})
      & \py{targets\_match\_helpers\_oracle} \\
    Identity, negative symmetry and composition hold on real batches
    (\cref{eq:inv-identity}-\cref{eq:inv-composition})
      & \py{invariants\_on\_batch} \\
    \addlinespace
    Quad rotation is an exact pixel permutation (\cref{eq:quad-k})
      & \py{crop\_patches\_quad\_is\_exact\_rot90} \\
    Batched extraction equals the per-patch loop (\cref{eq:inverse-warp})
      & \py{crop\_patches\_rotated\_} \\
      & \quad\py{batched\_equals\_loop} \\
    Angles are drawn from the stated sets (\cref{eq:angle-sampling})
      & \py{bounded\_rotation\_sampling} \\
    \addlinespace
    With rotation off the loss is unchanged (\cref{eq:loss-none})
      & \py{loss\_balance\_default\_unchanged} \\
    Variance balancing rescales both groups (\cref{eq:loss-variance})
      & \py{loss\_balance\_variance\_} \\
      & \quad\py{rescales\_groups} \\
    $\epsilon$ prevents the small-range blow-up (\cref{eq:loss-variance})
      & \py{loss\_eps\_guards\_small\_} \\
      & \quad\py{variance\_blowup} \\
    \addlinespace
    \py{ch2} controls rotate pixels but do not supervise them
    (\cref{sec:interpolation-controls})
      & \py{ss\_ch2\_control\_rotates\_} \\
      & \quad\py{but\_unsupervised} \\
    Metrics agree with a hand-computed oracle
    (\cref{sec:protocol-and-metrics})
      & \py{eval\_metrics\_oracle} \\
    \bottomrule
  \end{tabular}
\end{table}

\begin{table}[htbp]
  \centering
  \caption[Accelerator parity against CPU]{The same code on both accelerators, each
  measured against its own machine's CPU reference: ViT-B/32, batch 2, $M = 2048$, with
  the pair set held fixed so only arithmetic differs. Timing is one
  forward--backward--step, mean of ten.}
  \label{tab:device-parity}
  \begin{tabular}{l l l}
    \toprule
                                   & MPS (Apple silicon)   & CUDA (RTX 2080 Ti, \py{gpu30}) \\
    \midrule
    \py{torch}                     & 2.12.1                & 2.12.1+cu130 \\
    Max abs.\ deviation from CPU   & $6.4 \times 10^{-7}$  & $5.8 \times 10^{-7}$ \\
    Relative deviation             & $2.9 \times 10^{-6}$  & $2.6 \times 10^{-6}$ \\
    Loss vs CPU                    & identical (16 s.f.)   & rel.\ $1.1 \times 10^{-7}$ \\
    \py{randperm} matches CPU      & no                    & no \\
    \py{float64}                   & unsupported           & supported \\
    Seconds per step               & $0.110$               & $0.038$ \\
    \midrule
    Role                           & development, testing   & all reported pretraining \\
    \bottomrule
  \end{tabular}
\end{table}

\apxsection{Why the HLW AUC cannot carry the orientation claim}\label{apx:auc-weighting}
\Cref{sec:hlw-task} states that HLW's published AUC is structurally dominated by $\rho$ and
so cannot carry an orientation claim. This is the derivation. The horizon is the line $-x\sin\theta + y\cos\theta = \rho$, so on a square
crop ($W/H = 1$) its intercepts at the left and right image edges, in image heights, are
\begin{equation}\label{eq:hlw-lr}
  l \;=\; \frac{\rho - \tfrac{1}{2}\sin\theta}{\cos\theta}
    \;=\; \rho\sec\theta - \tfrac{1}{2}\tan\theta,
  \qquad
  r \;=\; \rho\sec\theta + \tfrac{1}{2}\tan\theta ,
\end{equation}
and the protocol's per-image error is the larger of the two edge discrepancies between the
predicted and true horizon, $e_{\texttt{horizon}} = \max\bigl(|\hat l - l|,\; |\hat r - r|\bigr)$,
from which the AUC is computed. Differentiating \cref{eq:hlw-lr} about $\theta = 0$ gives the
weight each channel carries into that error:
\begin{equation}\label{eq:hlw-weights}
  \Bigl|\frac{\partial l}{\partial \rho}\Bigr|_{\theta = 0} = 1,
  \qquad
  \Bigl|\frac{\partial l}{\partial \theta}\Bigr|_{\theta = 0} = \tfrac{1}{2} .
\end{equation}
An orientation error therefore enters the headline metric at \emph{half} weight and a
translation error at \emph{full} weight. Evaluated at the constant-predictor floor of the
held-out split ($\theta_0 = 1.643^\circ$, $\rho_0 = 0.3970$), the two contributions are
\begin{equation}\label{eq:hlw-dominance}
  \tfrac{1}{2}\tan\theta_0 \;=\; 0.0143
  \quad\text{against}\quad
  \rho_0 \;=\; 0.3970 ,
  \qquad
  \text{a ratio of } \approx 28\times .
\end{equation}
The published AUC is, in effect, a $\rho$-prediction score. Nothing in the derivation depends
on which arm is being scored or on any measured result, which is why the reporting choice
could be, and was, fixed before the four arms ran.


\end{document}
